\section{Related literature}\label{sec:literature}

\paragraph{Debt limits with endogenous rates.} \citet{mian2025} derive fiscal space when $R-G$ rises with debt. Their free-lunch condition, $R<G-\phi$ with $\phi$ the sensitivity of $R-G$ to log debt ($\psi b$ in our notation), is the $\varphi=0$ case of ours; their maturity extension works through the convenience yield of long versus short debt in steady state, not through repricing speed. \citet{lorenzoni2019} show that responsive fiscal rules and longer maturity can rule out slow-moving debt crises, through default equilibria. \citet{li2026} show that QE can shift a default boundary inward by depleting central-bank capital; they note that in transition longer maturity reduces rollover pressure, since only a fraction of the stock comes due each period, while in steady state it shrinks the fiscal dividend from convenience yields. Our setting has no default. Maturity leaves the stability threshold unchanged and instead governs the inflation alternative. In the tradition of \citet{sargent1981}, \citet{davig2011} and \citet{leeper2011}, the non-default limit is where fiscal adjustment stops and inflation must take over, and we make that boundary measurable. Related work on sustainability and debt capacity includes \citet{blanchard2019}, \citet{mehrotra2021}, \citet{reis2021} and \citet{ghosh2013}, and the convenience-yield literature \citep{krishnamurthy2012,choi2026,jiang2024}.

\paragraph{Consolidated maturity and QE.} \citet{greenwood2014} measure consolidated Treasury--Fed duration and show that the Treasury's maturity extension offset about a third of QE's reduction in privately held duration (35\% from a 2007 baseline; 63\% from December 2008). The \citet{tbac2020} computed consolidated maturity and duration, and since 2022 the Treasury's refunding materials report a consolidated weighted-average next rate reset, which \citet{tbac2026} reviews; it treats the part of the Fed's portfolio that funds reserves, reverse repos and other interest-bearing Fed liabilities as repricing overnight. These measures are scalar summaries (a weighted mean and, since 2024, a weighted median), shown as time series, and they exclude currency. The \citet{obr2021} documents the same shortening for the UK, and \citet{cbo2022}, \citet{levin2022}, \citet{cavallo2019} and \citet{davernas2024} discuss the fiscal risk it creates. \citet{delnegro2015} and \citet{hall2015} ask when a central bank's balance sheet needs fiscal support, and \citet{bassetto2013} trace the fiscal consequences of paying interest on reserves; our consolidated clock shows those consequences in the erosion base, since remunerated reserves took the Fed's book out of the inflation-tax base in 2008. In July 2022, Fed staff projected the deferred asset to peak at about \$60 billion, with a tail of about \$180 billion \citep{anderson2022}. We generalize the scalar reset measures to the full repricing profile, add currency to the erosion base, build the profile from 1980, and test it on realized outcomes, including the Fed's 2022--25 losses. On the monetary side, \citet{andreolli2024} shows that longer debt duration weakens the transmission of monetary policy to output in the United States and the UK, through a financing channel, and \citet{johns2026} find that the level and maturity of public debt shape transmission in Europe.

\paragraph{Maturity and the inflation tax.} In the fiscal theory of the price level, maturity determines whether fiscal news shows up as inflation now or later \citep{cochrane2001,cochrane2022,cochrane2023}, and monetary and fiscal policy jointly determine the price level \citep{sims2013}. \citet{barro2026} scale fiscal spending shocks by debt times duration in explaining 2020--23 inflation. \citet{hilscher2022} and \citet{aizenman2011} show that short U.S. maturity limits how far inflation can erode debt; \citet{hilscher2022} also show that financial repression through zero-interest reserves acts as a tax on their holders and makes inflation more effective. \citet{reis2017a} argues that QE, by bringing more debt due, lowers the price increase a fiscal crisis requires, and \citet{reis2017b} sets out how far a central bank can relieve fiscal burdens, including through seigniorage on the currency base. \citet{krause2016} show in a New Keynesian model with a maturity structure that a higher inflation target lowers real debt substantially only if the change is close to permanent; our inflation layer measures the same channel on the actual consolidated balance sheet, with currency in the erosion base. Appendix~\ref{app:whichinflation} reconciles the signs in this literature: a one-time level jump is maturity-neutral, a transitory surprise favors short debt, and sustained inflation favors long debt. We invert the erosion calculation into a required inflation rate as a function of the measured clock, add the monetary response, and quantify how consolidation shifts both.

\paragraph{Optimal maturity.} \citet{angeletos2002}, \citet{bhandari2017} and \citet{faraglia2019} study how the maturity structure should be chosen to insure the budget. We take the structure as given and measure its consequences.

\paragraph{Fiscal responses.} \citet{bohn1998,bohn2008} estimates a positive U.S. response of surpluses to debt; his coefficient maps exactly into ours (Appendix~\ref{app:fiscal}). \citet{eichengreen2026} find that surpluses respond to debt service more than to debt stocks, and more strongly when $r>g$; their object is a log-log elasticity, and they describe their evidence as descriptive. \citet{auerbach2024} find that legislated U.S. fiscal feedback, including the response to lagged net interest, was substantial before 2004 and has been about zero since. We use their estimates as a reference path, bring our own evidence in Appendix~\ref{app:fiscal}, and report every result as a function of the response.
